\documentclass{article}

\usepackage{amsmath}
\usepackage[boxed,vlined,boxruled]{algorithm2e}


\title{Range Queries}
\author{Kanishk Goel}
\date{1 June, 2026}
\parindent 0px

\SetKwInput{KwAssume}{Simplifying assumption}
\SetAlgoCaptionLayout{centerline}

\begin{document}
	\SetAlgoCaptionSeparator{}
	\renewcommand{\thealgocf}{} 
	\renewcommand{\AlCapNameFnt}{\bfseries}
	\renewcommand{\AlCapFnt}{}
	\DontPrintSemicolon
	\renewcommand{\algorithmcfname}{}
	\maketitle
	\SetAlgoNoEnd
	
	\section{Range Queries}
	The task is to calculate a value based on a subarray of an array. A naive way to answer a range query is to iterate over $i$ to $j$ such that $i < j$ and find the value that the query asks for. This runs in $O(n)$ and for $q$ queries will require $O(nq)$ time. Notation - $func_q(a_1,a_2 \ldots a_n)$
	
	\section{Sum Queries - Prefix Sum Subarray}
	Range query where value required is simply the sum of the elements in that range. This is solved in $O(n)$ via prefix sum array. \newline
	
	A variant of the same is to generalize the idea in higher dimensions. For example in $2D$, we can require the sum of a rectange from $(i,j)$ to $(i+l,j+b)$. \newline
	
	The solution to this, is to simply calculate a prefix sum subarray $f$, which calculates sum of the rectangle from $(0,0)$ to $(i,j)$. Now the answer would simply be $f(i,j) + f(i+l,j+b) - f(i,j+b) - f(i+l,j)$.
	
	\section{Min Queries - Sparse Table}
	This is a $O(n \, log \, n)$ preprocessing method, after which we can answer most queries in $O(log \, n)$ time and any minimum/maximum query in $O(1)$ time. \newline
	
	\subsection{Idea}
	Every number has a binary representation, therefore precalculate all values of $min_q(a,b)$ where $b-a+1$ is a power of $2$. \newline
	
	Every $2^k$ range can be calculated by using the previous $2^{k-1}$ ranges as 
	\[
	min_q(a,b) = \text{min}(min_q(a,a+w-1), min_q(a+w,b)),
	\]
	where $w = (b-a+1)/2$ \newline
	
	Now any query $(a,b)$ can be solved by representing the range $[a,b]$ as \newline $[a,a+k-1] \cup [b-k+1,b]$ where $k$ is the largest power of $2$ that does not exceed $b-a+1$.
	
	\[
	min_q(a,b) = \text{min}(min_q(a,a+k-1),min_q(b-k+1,b)),
	\]
	
	\subsection{Implementation}
	We use a $2D$ array, where $st[i][j]$ stores the answer for the range $[j,j+2^i - 1]$. Size of the array is $(K+1) \times MAXN$ where $MAXN$ is the biggest possible array length and $K \geq \lfloor log_2 \, MAXN \rfloor$. $MAXN$ dimension is second for cache friendliness.
	
	\subsection{Conclusion}
	Sparse Table is therefore, best used for idempotent functions (functions whose repeated application doesn't change the result beyond the first time).
	
	\section{Fenwick Tree}
	
\end{document}